% En esta seccion se realiza una investigación formal del tema y lo presenta bien fundamentado, haciendo referencia a las fuentes de información usando las normas APA.

\section{Introducción}
En esta segunda práctica del Laboratorio de Computación Gráfica el trabajo se movió del plano a las tres dimensiones. En la práctica anterior nos ocupamos de definir geometrías sencillas en coordenadas normalizadas de dispositivo; aquí el objetivo fue construir un objeto en el espacio y decidir desde dónde y cómo se le observa, lo cual introduce una separación que antes no existía: la geometría se define una sola vez y en su propio sistema de coordenadas, mientras que la imagen final depende de la cámara virtual y de las matrices que la describen. OpenGL es precisamente la interfaz que permite describir esos objetos y operaciones de forma independiente del equipo donde se ejecute el programa, delegando en el sistema el trabajo de convertirlos en píxeles \cite{grantmestrength_opengl}.

La cámara virtual queda determinada por tres vectores: su posición dentro de la escena, el punto hacia el que dirige la mirada y un vector de orientación que indica cuál dirección se despliega hacia la parte superior de la ventana. A partir de ellos se construye la matriz de vista, que lleva las coordenadas del mundo al sistema de referencia del observador. A esa matriz se suma la de proyección, encargada de delimitar el volumen de visualización y de aplanar la escena sobre el plano de despliegue, y que puede ser ortogonal, cuando los proyectores son paralelos y las medidas se conservan, o en perspectiva, cuando convergen en un punto y aparece el escorzo que asociamos con la visión humana. Finalmente, la matriz de modelo permite trasladar, rotar y escalar la geometría sin tocar sus vértices originales, de modo que un vértice llega a la pantalla después de la composición de las tres.

Cada vértice de la escena lleva asociado, además de su posición, un color expresado en el modo RGBA, en el que las tres componentes cromáticas y la de opacidad se manejan como valores flotantes acotados en lugar de las representaciones enteras habituales en otros contextos \cite{microsoft_rgba}. Por esa razón, los colores que elegimos a partir de códigos hexadecimales tuvieron que convertirse previamente al formato normalizado \cite{convertcolor_rgba}. Como el color se declara por vértice, la etapa de rasterización interpola los valores intermedios sin necesidad de código adicional, lo que produce los degradados que se aprecian en las capturas del reporte; en aplicaciones más elaboradas esa apariencia se resuelve mediante texturas, es decir, imágenes que se mapean sobre la superficie de la geometría \cite{uv_texturas}.

Las actividades de la práctica recorren precisamente esa cadena. Primero construimos la figura solicitada, una casa formada por un cuerpo prismático y un techo, cuidando que quedara centrada en el origen. Después configuramos las vistas frontal, lateral y superior asignándolas a las teclas F1, F2 y F3, y generamos variaciones de las proyecciones ortogonal y en perspectiva modificando sus parámetros. En seguida implementamos el movimiento de la cámara con las teclas WASD, y por último atendimos el modelo mismo: su desplazamiento sobre el eje Z con las flechas del teclado y su escalamiento con las teclas de más y menos. Para cada actividad se realizó además un experimento con el fin de observar cómo responde la escena al variar los valores involucrados.

El desarrollo se llevó a cabo en C++ sobre el proyecto base proporcionado por el profesor, que emplea OpenGL 3.3 con GLFW para la gestión de la ventana, GLAD para la carga de las funciones de la API y GLM para las operaciones con matrices y vectores. La compilación se realizó en Visual Studio siguiendo la configuración descrita en el manual de la práctica. Durante el desarrollo se utilizó además un modelo de lenguaje como apoyo para comprender la implementación del giro de la cámara y para la revisión de redacción del documento \cite{anthropic2026}.